% Chapter 2: Background and Prior Work (Unified)
% Combines: literature review + MOSD summary + MutClust summary
\chapter{배경 및 선행 연구}
\label{ch:background}

본 장에서는 본 학위논문의 핵심 연구인 MutBench를 이해하기 위한 배경을 네 개의 절로 구성하여 제시한다.

먼저, 다중오믹스 통합 분석(Section~\ref{sec:bg_moi})과 바이러스 변이 핫스팟 탐지(Section~\ref{sec:bg_hotspot})에 관한 선행 연구를 고찰한다.
이 두 분야는 본 학위논문의 핵심 연구(MutBench)와 선행 연구(MOSD, MutClust)가 다루는 핵심 영역이며, 공통적으로 클러스터링 기반 방법에 의존하면서도 체계적 벤치마킹이 부족하다는 문제를 공유한다.

이어서, 본 학위논문의 선행 연구로서 출판된 두 논문---MOSD (Section~\ref{sec:mosd_summary})와 MutClust (Section~\ref{sec:mutclust_summary})---를 요약한다.
MOSD는 다중오믹스 통합을 위한 체계적 벤치마킹 방법론을 확립하였고, MutClust는 바이러스 변이 핫스팟 탐지를 위한 적응형 알고리즘을 개발하였다.
본 장에서는 이 두 연구의 성과와 한계를 모두 제시하여, MutBench 개발의 동기와 방법론적 기반을 명확히 한다.


%% ===================================================================
%%  SECTION 1: 다중오믹스 통합 분석
%% ===================================================================
\section{다중오믹스 통합 분석}
\label{sec:bg_moi}

\label{subsec:bg_moi_methods}
\label{subsec:bg_benchmarking}

Multi-Omics Integration (MOI)은 유전자 발현, DNA 메틸화, 복제수 변이 등 다층적 생물학적 데이터를 결합하여 복잡한 질병의 분자적 특성을 포괄적으로 규명하기 위한 접근법이다.
그래프 기반(SNF~\cite{wang2014snf}, NEMO~\cite{rappoport2018nemo}), 행렬 분해 기반(MOFA~\cite{argelaguet2018mofa}, iClusterPlus~\cite{shen2013iclusterplus}), 합의 기반(COCA~\cite{hoadley2014coca}, PINSPlus~\cite{nguyen2017pinsplus}) 등 다양한 MOI 방법이 개발되었으나, 어떤 조건에서 어떤 방법이 최적인지에 대한 근본적 질문이 남아 있었으며, 기존 벤치마크 연구들~\cite{chauvel2020sample, rappoport2018benchmark, pierrejean2020feature, tini2019multi, duan2021evaluation}은 각각 1--3개의 요인만을 독립적으로 조사하여 통합적 가이드라인을 제공하지 못하였다.
저자는 이 공백을 해결하기 위해 MOSD 프레임워크~\cite{han2025mosd}를 개발하여, 10개 TCGA 암종에 걸쳐 10개 MOI 방법에 대한 9개 요인의 영향을 체계적으로 평가하였다.
MOSD는 두 가지 핵심 설계 원칙을 도입하였다: 세 가지 상보적 클러스터링 지표를 동등 가중치로 결합한 복합 평가 지표(cluster-score)와, 평가 대상 요인만 변동시키고 나머지를 통제하는 단일 요인 분리 설계이다.
Weber et al.~\cite{weber2019benchmarking}은 엄격한 벤치마킹에 중립적 평가, 다중 지표, 통제된 실험 설계, 재현성이 필요함을 제시하였으며, MOSD는 이러한 원칙을 MOI에 구현하였다. 이와 동일한 설계 원칙이 바이러스 핫스팟 탐지를 위한 MutBench로 이전되며, 이는 Section~\ref{sec:mosd_summary}에서 상술한다.


%% ===================================================================
%%  SECTION 2: 바이러스 변이 핫스팟 탐지
%% ===================================================================
\section{바이러스 변이 핫스팟 탐지}
\label{sec:bg_hotspot}

\subsection{변이 핫스팟 탐지 방법론}
\label{subsec:bg_hotspot_methods}

변이 핫스팟은 유전체 내에서 변이가 비균일적으로 집중되는 영역으로, 구조적 또는 면역학적 선택압 하에 있는 기능적으로 중요한 부위를 반영할 가능성이 높다.
바이러스 유전체에서 핫스팟을 식별하는 것은 변이체 출현 예측, 백신 표적 선정, 치료제 개발에 핵심적이다.
기존의 핫스팟 탐지 접근법은 빈도 기반, 엔트로피 기반, 밀도 기반 방법으로 분류할 수 있다.

\textbf{빈도 기반 접근법}은 아미노산 또는 뉴클레오타이드 수준에서 변이 빈도를 계산하고, 높은 빈도를 보이는 위치를 핫스팟으로 지정한다.
직관적이고 구현이 용이하지만, 두 가지 근본적 한계를 지닌다.
첫째, \textbf{계통 편향(lineage bias)}: 특정 계통이 데이터베이스에서 과대 대표될 경우 해당 계통의 정의 변이가 불균형적으로 높은 빈도 점수를 받는다.
D614G와 같이 거의 모든 현대 SARS-CoV-2 서열에서 관찰되는 변이는 실제 적응도 이점을 반영하지만, 근고정(near-fixation) 속도는 초기 팬데믹의 창시자 효과에 의해 증폭되어, 빈도 데이터만으로는 선택과 유전적 부동의 상대적 기여를 구분하기 어렵다.
둘째, \textbf{희귀 변이의 간과}: 낮은 빈도로 발생하지만 다양한 형태로 나타나는 변이---면역 회피 또는 바이러스 적응을 시사할 수 있는---가 체계적으로 간과된다.

\textbf{엔트로피 기반 접근법}은 Shannon entropy를 사용하여 주어진 위치에서 뉴클레오타이드 다양성을 측정한다.
Mullick et al.~\cite{mullick2021entropy}은 Shannon entropy와 K-means clustering을 결합하여 SARS-CoV-2 Spike 단백질의 핫스팟을 탐지하였다.
그러나 기존의 뉴클레오타이드 엔트로피는 변이 발생 여부와 관계없이 모든 뉴클레오타이드의 분포를 반영하므로, 변이 발생을 조건으로 한 실효적 다양성을 포착하지 못한다.
또한 K-means clustering은 클러스터 수를 사전에 지정해야 하며 밀도 구조를 고려하지 않는다.

\textbf{암 유전체학 방법.}
암 유전체학에서는 체세포 변이의 핫스팟을 탐지하기 위한 다양한 방법이 개발되었다.
MutSigCV~\cite{lawrence2013mutsigcv}는 배경 변이율을 모델링하여 과도한 변이를 보이는 유전자를 식별하나, 유전자 수준의 해상도에서 작동한다.
OncodriveCLUST~\cite{tamborero2013oncodriveclust}는 변이의 위치적 클러스터링 편향을 무작위 기대값과 비교한다.
HotMAPS~\cite{tokheim2016hotmaps}는 변이를 3D 단백질 구조에 매핑하여 공간적 핫스팟을 탐지하나, 단백질 구조 데이터가 필요하며 코딩 영역에 제한된다.
MutSpot~\cite{munoz2020mutspot}은 후성유전체적 특성을 활용하는 random forest 모델로 변이 밀도를 예측한다.
이러한 암 유전체학 방법들은 바이러스 유전체에 직접 적용하기 어렵다: 바이러스 유전체는 암 유전체보다 1,000배 이상 작으며, 변이 밀도 패턴, 선택압의 유형, 진화 동역학이 근본적으로 다르다.

\subsection{밀도 기반 클러스터링}
\label{subsec:bg_density_clustering}

DBSCAN~\cite{ester1996dbscan}은 고밀도 영역을 클러스터로 정의하고 저밀도 영역을 노이즈로 분류하는 대표적인 밀도 기반 클러스터링 알고리즘이다.
두 가지 핵심 파라미터---epsilon(탐색 반경)과 MinPts(최소 포인트 수)---를 사용하여 밀도 조건을 정의한다.
그러나 DBSCAN은 \textbf{고정된 전역 epsilon}을 사용하므로, 밀도 변동이 큰 데이터에서는 적절한 클러스터를 형성할 수 없다.
이 문제는 바이러스 유전체에서 극단적인 형태로 나타난다:
SARS-CoV-2 유전체에 직접 적용하면 데이터의 99.76\%가 단일 거대 클러스터에 속하는 총 4개의 클러스터만 형성되어, 생물학적으로 의미 있는 핫스팟 구별이 불가능해진다~\cite{youn2025mutclust}.

HDBSCAN~\cite{campello2013hdbscan}은 DBSCAN의 고정 epsilon이라는 근본적 한계를 해결한다.
다양한 epsilon 값에 걸쳐 DBSCAN 결과를 계층적으로 탐색하여 밀도 변동에 적응적으로 클러스터를 추출한다.
단일 파라미터인 min\_cluster\_size로 작동하며, epsilon의 명시적 지정이 필요하지 않다.
임의 형태의 클러스터를 탐지할 수 있고, 노이즈 포인트를 자동으로 식별하며, 클러스터 소속의 확률적 해석을 제공한다.

OPTICS (Ordering Points To Identify the Clustering Structure)~\cite{ankerst1999optics}는 도달 가능성 플롯(reachability plot)을 생성하여 모든 밀도 수준의 클러스터링 구조를 동시에 드러냄으로써 DBSCAN의 한계를 해결한다.
밀도 계층으로부터 평탄 클러스터를 자동 추출하는 HDBSCAN과 달리, OPTICS는 클러스터 추출을 위해 후처리(예: $\xi$-steepness 방법)가 필요하여 자동화된 파이프라인에서의 활용이 제한적이다.

빈도 기반 변이 분석의 중요한 한계는 \textbf{계통발생학적 비독립성(phylogenetic non-independence)}이다: 근연 서열에서 관찰되는 변이는 독립적 사건이 아니다~\cite{felsenstein1985phylogenies}.
TreeTime~\cite{sagulenko2018treetime}은 계통수 상에서 조상 서열을 재구성하여 독립적 변이 발생을 계수하며,
BEAST~\cite{suchard2018beast}는 진화 과정을 직접 모델링하는 베이지안 프레임워크를 제공하고,
HyPhy~\cite{kosakovsky2020hyphy}는 개별 코돈에서 선택압을 검정할 때 계통발생 구조를 고려한다.
이러한 계통발생 인식 방법들은 밀도 기반 핫스팟 탐지와 상호 보완적이나, 고품질 계통수를 필요로 하며 밀집 표본 데이터셋에서는 신뢰성이 제한적이다~\cite{turakhia2020ultrafast}.

\textbf{Deep Mutational Scanning (DMS)}은 단백질 내 가능한 모든 단일 아미노산 치환의 기능적 효과를 체계적으로 측정하는 실험 기법이다~\cite{fowler2014dms}.
SARS-CoV-2의 경우 Starr et al.~\cite{starr2020dms}이 수용체 결합 도메인(RBD)에서 ACE2 결합 친화도와 발현을 측정하였고, Dadonaite et al.~\cite{dadonaite2023dms}이 전체 길이 Spike 단백질에서 슈도바이러스 세포 진입 적응도를 측정하였다.
Bloom and Neher~\cite{bloom2023fitness}는 모든 SARS-CoV-2 단백질에 걸친 DMS 적응도 효과의 통합 분석을 제공하였다.
DMS 데이터는 핫스팟 벤치마킹을 위한 자연스러운 ground truth를 제공한다: 변이가 큰 적응도 효과를 보이는 위치는 탐지 방법이 식별해야 하는 생물학적 기능 부위에 해당한다.
이는 다중오믹스 벤치마킹에서 분자적으로 정의된 암 아형을 ground truth로 사용하는 원칙과 유사하다.

바이러스 유전체학에서 핫스팟 탐지 방법의 벤치마킹이 부재한 현황은 암 유전체학과 대조적이다.
암 유전체학에서는 Bailey et al.~\cite{bailey2018comprehensive}의 포괄적 드라이버 유전자 분석이 표준 벤치마크로 기능하며, 다수의 탐지 방법이 동일한 데이터셋과 ground truth에 대해 체계적으로 비교되어 왔다.
바이러스 유전체에 대해서는 그러한 벤치마크가 존재하지 않으며, MutClust를 포함한 방법들은 자기 참조적 통계 검정에만 의존한다.
이 공백이 본 학위논문의 핵심 연구인 MutBench (Chapter~\ref{ch:mutbench})의 근본적 동기이다.



%% ===================================================================
%%  SECTION 3: MOSD 요약
%% ===================================================================
\section{MOSD: 다중오믹스 연구 설계 프레임워크}
\label{sec:mosd_summary}

\textit{본 절은 BMC Genomics (2025) 26:769~\cite{han2025mosd}에 출판된 논문의 요약이다.}

\label{subsec:mosd_bg}
\label{subsec:mosd_methods_summary}
\label{subsec:mosd_results_summary}
\label{subsec:mosd_implications}

MOSD (Multi-Omics Study Design)~\cite{han2025mosd}는 10개 TCGA 암종에 걸쳐 10개 MOI 방법에 대한 9개 요인의 영향을 체계적으로 평가한 최초의 포괄적 벤치마킹 프레임워크이다.
핵심 기여는 세 가지이다:
(1)~세 가지 상보적 클러스터링 지표(ARI, NMI, F-measure)를 동등 가중치로 결합한 \textbf{복합 cluster-score};
(2)~평가 대상 요인만 변동시키고 나머지를 통제하는 \textbf{단일 요인 분리 설계};
(3)~\textbf{연구 설계가 알고리즘 선택보다 중요하다}는 핵심 발견---어떤 단일 방법도 모든 조건에서 우위를 보이지 않아, 방법 선택이 본질적으로 맥락 의존적임을 확인하였다.

본 학위논문에서 MOSD의 주된 기여는 \textbf{복합 지표 설계 원칙}이다: cluster-score의 상보적 지표 동등 가중 결합이 MutBench의 hotspot-score $= (\text{recall} + \text{precision} + \text{stability}) / 3$으로 직접 이전된다.
MOSD의 단일 요인 분리 설계 역시 MutBench의 파라미터 민감도 분석에 적용된다.
평가 요인을 체계적으로 통제하고 다중 지표로 평가해야 한다는 프레임워크 수준의 원칙이 MutBench의 설계를 직접적으로 형성한다.


%% ===================================================================
%%  SECTION 4: MutClust 요약
%% ===================================================================
\section{MutClust: 적응형 밀도 기반 핫스팟 탐지}
\label{sec:mutclust_summary}

\textit{본 절은 BioData Mining (2025) 18:61~\cite{youn2025mutclust}에 출판된 논문의 요약이다.}

\subsection{H-score와 적응형 DBSCAN}
\label{subsec:mutclust_hscore}

MutClust는 바이러스 유전체의 변이 핫스팟 탐지를 위해 설계된 최초의 밀도 기반 알고리즘이다.
핵심 혁신은 새로운 변이 중요도 지표인 H-score와, 이에 기반한 적응형 DBSCAN 프레임워크이다.

\textbf{H-score}는 변이 빈도와 변이 엔트로피(변이 발생을 조건으로 한 다양성)의 곱에 로그 변환을 적용하여 두 가지 상보적 신호를 동시에 포착하며, 높은 빈도이지만 낮은 다양성의 계통 정의 변이와, 낮은 빈도이지만 허위적으로 높은 다양성의 노이즈 위치라는 두 유형의 위양성을 억제한다.

H-score에 기반한 \textbf{적응형 DBSCAN}은 DBSCAN의 고정된 전역 epsilon을 각 위치의 H-score에 비례하는 이웃 반경으로 대체하여, 생물학적으로 중요한 영역에는 넓은 탐색 반경을, 그 외 영역에는 좁은 반경을 할당한다.
알고리즘은 먼저 개별 H-score, 이웃 평균, 이웃 합, 최소 포인트 수에 대한 임계 조건을 충족하는 Core Candidate Mutation (CCM)을 식별한 후, 각 CCM으로부터 축소되는 반경으로 양방향 클러스터 확장을 수행한다.
핵심 파라미터($\gamma=10$, $d=3$, MinPts$=5$)는 224,318개 SARS-CoV-2 서열로부터 경험적으로 결정되었으며, 전체 시간 복잡도는 $O(N + M \log M)$으로 $N$은 유전체 길이, $M$은 변이 위치 수이다.

\subsection{SARS-CoV-2 핫스팟 탐지 결과}
\label{subsec:mutclust_results_summary}

224,318개 GISAID SARS-CoV-2 서열을 우한-Hu-1 참조 유전체(NC\_045512.2, 29,903 bp)에 정렬하여 분석한 결과, MutClust는 \textbf{477개}의 변이 핫스팟을 탐지하였으며, 이는 DBSCAN으로 얻은 4개 클러스터와 질적으로 완전히 다른 수준의 해상도이다.
검증을 위해 조사한 10개 알려진 기능적 변이 중 D614G, E484K, N501Y, L452R 등 9개가 핫스팟에 포함되었으며, Spike 단백질 영역에서 가장 높은 핫스팟 밀도가 관찰되었다.
인플루엔자에 대한 검증(276,910개 서열, 3개 아형)에서도 D222G(H1N1), E627K(H3N2), K338R(B-Victoria) 등 임상적으로 중요한 변이가 핫스팟으로 성공적으로 탐지되어 일반화 가능성이 확인되었다.
특히 H3N2의 경우, Smith et al.~\cite{smith2004antigenic}이 항원 지도작성(antigenic cartography)을 통해 HA 단백질의 소수 치환이 항원 변화를 주도함을 밝힌 바 있어, 이러한 핵심 위치의 정확한 탐지가 중요하다.
그러나 인플루엔자 유전체는 SARS-CoV-2의 약 10분의 1 크기이므로, 최적 성능을 위해서는 유전체 특성에 맞춘 파라미터 재조정이 필요할 수 있다.

\subsection{한계점과 MutBench 개발 동기}
\label{subsec:mutclust_limitations}

MutClust의 성과에도 불구하고, 세 가지 근본적 한계가 MutBench 개발의 직접적 동기가 된다.

\textbf{첫째, 구현 수준의 문제이다.}
클러스터 확장 함수의 diminishing 공식에 존재하는 버그로 인해 선형이 아닌 누적적 반경 감쇠가 발생하여 클러스터 확장이 조기에 종료되며, 특정 경계 조건에서는 감쇠 방향이 역전되는 현상이 나타난다.

\textbf{둘째, 하드코딩된 파라미터이다.}
$\gamma=10$, $d=3$, CCM 임계값(0.03, 0.01, 0.05)은 224,318개 SARS-CoV-2 서열에 맞춰 경험적으로 설정된 값이다.
서열 수나 유전체 크기가 달라지면 이 값들이 부적절해질 수 있으나, 파라미터 변동에 따른 탐지 결과의 민감도가 체계적으로 분석된 적이 없다.
인플루엔자 검증에서 동일 파라미터의 적용이 성공적이었지만, 인플루엔자 유전체는 SARS-CoV-2의 약 10분의 1 크기($\sim$2,300 bp/분절 대 $\sim$30,000 bp)이므로 최적 성능을 위한 파라미터 재조정의 필요성이 불확실한 채로 남아 있다.

\textbf{셋째, 체계적 벤치마킹의 부재이다.}
MutClust의 핫스팟 탐지 성능은 알려진 기능적 변이와의 정성적 비교에만 의존하며, 독립적 ground truth에 기반한 정량적 평가(정밀도, 재현율, 안정성)가 수행되지 않았다.
암 유전체학에서는 Bailey et al.~\cite{bailey2018comprehensive}의 포괄적 드라이버 유전자 분석과 같이 핫스팟 탐지 방법의 체계적 비교 연구가 수행되었으나, 바이러스 유전체에 대해서는 그러한 벤치마크가 존재하지 않는다.
MutClust를 포함한 바이러스 핫스팟 탐지 방법들은 자기 참조적 통계 검정(bootstrap)에만 의존하며, 표준화된 평가 지표의 부재로 인해 서로 다른 방법의 결과를 공정하게 비교하는 것이 불가능하다.

추가적으로, H-score가 뉴클레오타이드 수준에서 산출되므로 동의 변이와 비동의 변이를 구분하지 않으며, dN/dS 비율과 같은 선택압 지표와의 통합이 필요하다.
계통발생학적 비독립성---변이 빈도가 독립적 양성 선택이 아닌 단일 창시자 효과의 결과일 수 있음---도 보정되지 않는다.

\vspace{1em}

이러한 한계의 핵심은 MutClust가 ``새로운 탐지 방법''을 제시하면서도 ``그 방법의 성능을 정량적으로 평가할 프레임워크''를 함께 제시하지 못했다는 점이다.
이는 Section~\ref{subsec:bg_benchmarking}에서 고찰한 벤치마킹 원칙---중립적 평가, 다중 지표, 통제된 실험 설계---의 부재에 해당한다.

Chapter~\ref{ch:mutbench}의 MutBench는 구현 개선, 독립적 ground truth 구축, 복합 평가 체계, 파라미터 민감도 분석을 통해 이 세 가지 한계를 체계적으로 해결한다.
